$f_a'(x)=3(x^2-a)$。
当 $a<0$ 时，$f_a'(x)>0$，故在 $\mathbb R$ 上严格递增，无局部极值点。
当 $a=0$ 时，$f_0(x)=x^3+1$ 在 $\mathbb R$ 上严格递增，$x=0$ 虽为驻点但不是局部极值点。
当 $a>0$ 时，增区间为 $(-\infty,-\sqrt a]$、$[\sqrt a,+\infty)$，减区间为 $[-\sqrt a,\sqrt a]$，故恰有两个局部极值点，分别为 $-\sqrt a$、$\sqrt a$。
综上，$\boxed{a\le0\text{ 时为 }0\text{ 个；}a>0\text{ 时为 }2\text{ 个}}$。\qed
